\begin{proposition}[7.1.11]
\label{IV.2.s7.1.11-fr}
\begin{enumerate}
\item[\textup{(i)}] Todo anillo local noetheriano formalmente
catenario es universalmente catenario.
\item[\textup{(ii)}] Si $A$ es un anillo local noetheriano
formalmente catenario, todo anillo local $B$ que
sea una $A$-álgebra esencialmente de tipo finito
es formalmente catenario.
\end{enumerate}
\end{proposition}

\begin{proof}
La afirmación (ii) se deduce inmediatamente de la
condición \emph{c)} de (7.1.8) y del hecho de que, si un
anillo local $C$ es una $B$-álgebra esencialmente
de tipo finito, $C$ también es una $A$-álgebra
esencialmente de tipo finito (1.3.9).
Para demostrar (i), observemos primero que, por
la condición \emph{b)} de (7.1.8), un anillo local
noetheriano $A$ formalmente catenario es catenario,
pues los cocientes de $A$ por sus ideales primos
minimales lo son (7.1.4). Si ahora $E$ es una
$A$-álgebra de tipo finito, de lo anterior y de
(ii) se deduce que todo anillo local de $E$ en
un ideal primo es catenario, y por tanto $E$ es
catenario. Así pues, $A$ es universalmente
catenario.
\end{proof}
